\documentclass{article}
\usepackage[utf8]{inputenc}
\usepackage{enumitem}
\title{Testprosedyre for EMP-generator}
\author{}
\date{}

\begin{document}
\maketitle

\section{Sikkerhet}
\begin{itemize}
\item Bruk alltid vernebriller og isolerende hansker.
\item Ha en brannslukker tilgjengelig.
\item Utlad alltid kondensatoren før du berører kretsen.
\item Utfør tester i et kontrollert miljø, helst utendørs eller i et Faraday-bur.
\end{itemize}

\section{Funksjonstest av lader (hovedversjon)}
\begin{enumerate}
\item Koble fra gnistgap og spole, men la hovedkondensator C1 være tilkoblet.
\item Koble et voltmeter (1000 V DC) mellom utgang (D1 katode) og jord.
\item Slå på S3, hold S1. Spenningen skal stige til 400–450 V i løpet av 5–10 sek. LED skal lyse.
\item Slipp S1; spenningen skal holde seg. Trykk S2 og hold i minst 5 sekunder; spenningen skal falle til under 10 V.
\end{enumerate}

\section{Test av gnistgap}
\begin{enumerate}
\item Koble til hovedkondensator og gnistgap, men la spolen være frakoblet (åpen krets).
\item Lad til 400 V. Trykk på piezo-knappen. Du skal høre et skarpt smell og se en gnist.
\item Hvis ikke, juster gapet mindre eller sjekk piezo-tilkobling.
\end{enumerate}

\section{Test med spole}
\begin{enumerate}
\item Koble spolen til gnistgapets utgang.
\item Plasser en målespole (5 vindinger, 2 cm diameter) 20 cm foran.
\item Koble målespolen til oscilloskop (1 M$\Omega$, 10x probe).
\item Lad og fyr. Du skal se en dempet sinus med amplitude 10–50 V.
\item Mål frekvens og toppspenning. Frekvensen bør være rundt 3.6 kHz.
\end{enumerate}

\section{Ødeleggelsestest}
\begin{enumerate}
\item Plasser en offer-enhet (f.eks. Arduino, kalkulator) 5–10 cm fra spolen.
\item Lad og fyr. Enheten skal slutte å fungere.
\item Gjenta med økende avstand for å finne effektiv rekkevidde.
\end{enumerate}

\section{Måling av toppstrøm (valgfritt)}
\begin{enumerate}
\item Bruk en Rogowski-spole eller en strømtransformator rundt en av ledningene til spolen.
\item Koble til oscilloskop og integrer signalet for å finne strømmen.
\item Sammenlign med teoretisk verdi.
\end{enumerate}
\end{document}
